%-------------------------------------------------------------------------------
%	SECTION TITLE
%-------------------------------------------------------------------------------
\cvsection{Research Experience}


%-------------------------------------------------------------------------------
%	CONTENT
%-------------------------------------------------------------------------------
\begin{cventries}

%---------------------------------------------------------
\cventry
{}{Search for low-MET electroweak supersymmetry}{Texas Tech University}{2022-ongoing}{
\begin{cvitems}
\item {Main topic of the thesis: This analysis focusses on all-hadronic signatures of stealth and RPV decays of charginos and neutralino of the supersymmetry}
\item {The analysis will use CMS collision data from Run 2 (year 2016-2018) and Run 3 (year 2022-onwards)}
\item {These extremely hard to probe models gain sensitivity using novel deep learning jet taggers, increased luminosity and better understanding of the detector}
\item {Events with high energy gluon jets and top pairs are the dominant standard model backgrounds. A combination of particlenet bb, W jet taggers is used to improve sensitivity}
\item {This analysis searches for supersymmetry in previously unexplored region using cutting edge AI and machine learning techniques. If we succeed in observing SUSY signatures, it would represent a major breakthrough in particle physics. If we fail to find any signatures, we will be able to place stronger upper limits on its production cross section.}
\end{cvitems}
}

%---------------------------------------------------------
  \cventry
    {Research Coordinator} % Job title
    {Jnana Prabodhini Foundation} % Organization
    {Pune, India \& Lubbock, USA} % Location
    {2021} % Date(s)
    {
      \begin{cvitems} % Description(s) of tasks/responsibilities
        \item {A 3-year long project aimed at estimating COVID-19 mortality and undercount factor for Pune city in India.}
        \item {A frugal method of estimating COVID-19 mortality was tested against traditional statistics-based models widely used in epidemiology and media reports}
        \item {The data was collected through public surveying in 2 rounds over a period of a year}
        \item {The results were published in nature - scientific reports}
        \item {Gained experience in leading a team of researchers and working on topic outside the domain expertise.}
      \end{cvitems}
    }
\end{cventries}

%% \cvsection{Technical/Service Work}
%% \begin{cventries}
%% %---------------------------------------------------------
%%   \cventry
%%     {Prompt Feedback and Troubleshooting} % Job title
%%     {Hadronic Calorimeter (HCAL)} % Organization
%%     {CERN, Geneva} % Location
%%     {2024} % Date(s)
%%     {
%%       \begin{cvitems} % Description(s) of tasks/responsibilities
%%         \item {Member of an international group responsible for monitoring and troubleshooting HCAL operations}
%%         \item {Among the first ones to look at the data recorded by HCAL since 2023}
%%         %% \item {Helped identify and debug issues with HCAL on short time scales. Eg, lumi spikes in the forward part of HCAL, anomalous peak in total collected charge, effects of response corrections etc.}
%%         %% \item {Monitored and compiled daily pedestal (noise) trends. The values calculated were later used in downstream event reconstruction.}
%%         %% \item {Developed framework for analyzing HCAL raw data which is now part of HCAL analysis workflow.}
%%         \item {Identified and mitigated HCAL operational anomalies: Diagnosed high-frequency "lumi spikes" in the forward HCAL and investigated anomalous peaks in total collected charge to ensure data quality for Run 3.}
%%         \item {Managed detector stability via noise monitoring : Developed and maintained the daily pedestal (noise) trend monitoring framework, directly providing the calibration constants required for downstream event reconstruction.}
%%         \item {Authored HCAL Raw Data Analysis documentation : Developed a comprehensive analysis framework and documentation for HCAL raw data that has been formally adopted into the CMS collaboration’s internal literature.}
%%       \end{cvitems}
%%     }

%% %---------------------------------------------------------
%%   \cventry
%%     {Performance studies} % Job title
%%     {} % Organization
%%     {} % Location
%%     {2023} % Date(s)
%%     {
%%       \begin{cvitems} % Description(s) of tasks/responsibilities
%%         \item {Member of an international group dedicated to study the performance of HCAL of CMS}
%%         \item {Worked on the data collected in 2023 phase scan}
%%         \item {Parametrized and characterized the raw pulse shape in different parts of the detector}
%%         \item {The pulse shape parametrization is now part of the official CMS simulation framework}
%%       \end{cvitems}
%%     }

%% %---------------------------------------------------------
%%   \cventry
%%     {Test beam participant} % Job title
%%     {Light Dark Matter Experiment (LDMX)} % Organization
%%     {CERN, Geneva} % Location
%%     {2022} % Date(s)
%%     {
%%       \begin{cvitems} % Description(s) of tasks/responsibilities
%%         \item {Developed slow control software for the prototype trigger scintillator for the \href{https://confluence.slac.stanford.edu/display/MME/Light+Dark+Matter+Experiment}{LDMX}}
%%         \item {Helped Construct a cosmic ray test stand at Texas Tech University, recorded and analyzed data}
%%         \item {Participated in test beam at CERN, where different detector components of LDMX were tested using high energy pion, electron and muon beams}
%%         \item {Monitored and controlled collision runs in shifts.}
%%         \item {Assumed charge of one of the subdetectors called trigger scintillator}
%%       \end{cvitems}
%%     }
%% \end{cventries}
